\chapter{绪论}
\label{ch:intro}

\section{一体机收窄了哪些可测量}

蒸箱把饱和蒸汽送进封闭腔体。蒸汽在食品表面冷凝，潜热在那里释放，热量再靠导热由外向内推进。烤箱走另一条路：干热空气与管壁辐射把表面迅速推到褐变温度，内部却仍按厚度平方的时间尺度缓慢跟上。这两条边界本来分属两台机器，厨房台面因此被占去两处，蒸制与烤制也各自保留一套互不引用的定时曲线。蒸的时候水汽大、时间难校准，中心未熟与表面过湿交替出现；烤的时候腔温波动、搁架和门缝让边界失去对称，冷点可以偏离几何中心\cite{feyissa2013}。漏电、干烧和超温则是两类机器共有的隐患，现有保护回路读水位、门开关、剩余电流和加热器温度，与食品内部的进度估计分属两套互不相干的回路。

家用蒸箱的容积多在 \(20\)--\(30\,\mathrm{L}\)，饱和蒸汽由底部水箱或外置发生器供给，冷凝热流密度可达每平方米数千瓦。薄件的表面很快接近饱和温度，内部却按导热时间尺度慢慢跟上。对流烤箱把表面送到 \(150^\circ\mathrm{C}\) 以上，褐变和结壳只发生在近表面几毫米，冷点位置仍由厚度平方决定。两台机器各自标定一套菜单：蒸箱按重量或按“鱼、蔬菜、馒头”给分钟数，烤箱按腔温加定时。菜单一旦写进面板，这一次进门的初温、脂肪夹层和真实厚度就不再改写停机时刻。操作者看得见门缝逸出的蒸汽、发红的管壁和升高的风机噪声，读不到的是冷点有没有越过 \(70^\circ\mathrm{C}\)、表面含水有没有落到结壳门槛以下。

把蒸与烤并进同一块钢板，省下的是台面，收窄的是旋钮。外壳换成深色搪瓷或拉丝不锈钢，夹层塞岩棉，内衬做圆角；门玻璃内外两层，中间留一条窄气隙。顶部或背部是加热管，一侧是蒸汽喷嘴，后壁是循环风机，右上或顶部是排汽百叶。水箱藏在底部或侧舱，容量约 \(1\,\mathrm{L}\)，水位开关接入安全回路；门铰链带微动开关，门一开，加热与蒸汽同时归零。这些零件在厨房里都看得见。看不见的是食品沿厚度的温度剖面。操作者拧的是档位、温度和时间，机器真正改的只有四个量：加热功率、蒸汽阀开度、排汽开度和风机转速。四个执行量进门之后变成表面通量，冷点进度由厚度平方和这一次过程里看见的热阻共同决定。

于是问题落在一个具体的量上。食品一旦进门，冷点温度、内部含水和反应进度都不在四个执行量上；现有商品路径——机械定时、腔温闭环、按重量调用曲线、穿刺探针、表面红外——闭合的分别是空气温度、申报质量或刺入点的读数（止点见第\ref{ch:lit}章），而这一块食品的内部热阻，在这一次过程里从未进入后验。更麻烦的是符号冲突：蒸汽把表面打湿，加热把表面推干，同一只阀门对冷点和表面品质的作用方向可以相反。为冷点多送一股蒸汽，表面可能湿到无法结壳；为表面排湿而加大排汽，冷点又可能失去冷凝潜热。分体时代，这个矛盾被两台机器分摊；并进同一腔体之后，它必须在同一套执行量上解决。蒸档夹住加热，烤档关掉蒸汽，组合档把四个量一起放开：档位名称写在面板上，允许集合写在执行器接口上，食品模块收到的始终是通量。

家用过程还有一条时间上限。\(55\,\mathrm{min}\) 既是常见蒸烤时域，也是数值核的积分终点。操作者很少把一体机开过一小时，水箱容量、门封和电气温升都按这个量级设计。薄件的导热时间尺度在开环菜单的时间内已经走完，厚件的 Fourier 数在同一上限内还远小于 \(1\)。同样四个执行量，薄件被菜单推过设计线，厚件停在设计线之下。厚度把整条响应带平移，策略名称只是跟在厚度后面的标签。

比较因此被固定在同一个积分器、同一组四个执行量和同一条设计线上，也就是 \(\alpha_m=0.85\) 所定义的那条线；开环、腔温比例、无探针与对偶这四条策略都共用它。后面出现的交叉表、网格加密、因子消融和典型轨迹，比的都是这四条策略在十六条控制体上的停机偏差，而不是各自换一套口径。同一条状态流最终会被读回三维腔体，地址为 \url{https://2711944586.github.io/xxgz/}。

\section{内部热阻缺席带来的决策困难}

记食品与腔体的真实状态为 \(\vect{x}(t)\)，可安装传感器给出
\[
\vect{y}(t)=h(\vect{x}(t))+\vect{v}(t).
\]
停机、表面品质和安全构成决策泛函 \(\phi(\vect{x})\)。沿 Lindley 的实验信息定义\cite{lindley1956}，信息赤字为
\begin{equation}
\label{eq:deficit}
\Delta\mathcal{I}(t)=H\big(\phi(\vect{x}(t))\mid \vect{y}_{0:t}\big)-H_{\mathrm{tol}}.
\end{equation}
\(H_{\mathrm{tol}}\) 由动作切换的滞回带决定。\(\Delta\mathcal{I}>0\) 时若仍按开环时间表改边界，冷点进度和表面干耗会分开漂移：二者对同一蒸汽指令的符号不必相同。一块薄面团在饱和蒸汽里很快被表面冷凝推过设计线，同一股蒸汽打在厚肉糜上，冷点几乎不动，表面却已经被打湿。赤字落在当前观测对决策泛函的条件熵上：滞回带还没有被这一次过程看见的内部热阻填平。加密同一路腔温，只是把已经闭合的空气温度再画密一档。

机器要完成的链是
\begin{equation}
\label{eq:chain}
\begin{aligned}
&\text{边界激励}\to\text{多尺度状态}\to\text{后验}\\
&\to\text{信息与品质价值}\\
&\to\text{安全投影}\to\text{新的边界}.
\end{aligned}
\end{equation}
链上缺一环，后面的环就没有输入。没有后验，信息价值和品质价值无法分开；没有安全投影，品质模型的错误会直接变成执行器指令。现有商品路径多数在第一条就停住：它们改边界，却不把这一次过程里的内部热阻写进后验。

图\ref{fig:cavity}给出同一腔体上的四个执行量，加热器在顶部，蒸汽从一侧进入，风机带动气流，排汽带走湿空气。档位只限制哪几个量允许离开零：蒸档把加热请求夹在 \(0.25\) 以下，烤档把蒸汽固定在零并把排汽抬到至少 \(0.45\)，组合档放开四个量。水位、门、剩余电流和加热器温度是硬测量：其中任何一条越过阈值，加热和蒸汽同时被固定在零。进度、热阻和估计冷点留在品质回路里，硬切断的自变量只有这四条测量。厚件在家用时域内走不完导热时间尺度：探针可以把热阻后验收窄，扩散时间仍把厚件压在设计线之下。

\begin{figure}[htbp]
\centering
\colimg{S_cavity.png}
\caption{同一腔体上的四个执行量。食品沿厚度离散为十六条控制体，冷点取最低温度}
\label{fig:cavity}
\end{figure}

\section{从边界激励到安全投影的五层结构}
\label{ch:arch}

四个执行量是闭环唯一的出口，但进入闭环的两路信号必须分开：表面温度、腔温和质量从食品侧进来，水位、管温、门位、腔压与剩余电流则从安全侧进来。两侧不共用同一根信号线，品质模型的错误改写不了断电条件，断电条件也回写不了进度后验。家电面板上那些被称为“智能”的功能，多数恰恰把这两侧混写进同一条加热请求，结果超温切断会跟着后验一起抖动。品质回路只该改请求，安全回路只该改指令。

在产品上，这件事对应两块彼此独立的电路。一块读 NTC 或热电偶给出的加热器温度、水箱浮子给出的水位、门微动开关与漏电互感器，任一越过阈值就切断加热管和蒸汽阀；另一块读表面温度或腔温，用来判断还要不要继续加热。两块若共用同一条请求线，麻烦就来了：进度估计一旦偏高，加热器温度保护也会跟着松；估计一旦偏低，超温切断又可能把已经接近设计线的过程掐断。已经运行的核因此把请求与指令分成两列，安全投影插在中间，自变量只有硬测量。

图\ref{fig:loop}给出依存关系，食品与腔体边界产生可安装观测，观测更新决策状态后验；后验一并进入对偶决策和安全投影。信息动作与品质动作在投影之前分开计算，进入同一个执行器接口后写回边界。已经运行的核可以压缩成五段：边界、十六条控制体、表面观测、探针或跟踪、加热指令，加热指令再写回边界；结构方程的路径积分不在这五段里，这条顺序由下一段的时间步直接给出。图\ref{fig:chain}把这条链本身画成六环，并标出环上每一环由哪一条设计性质约束。

\begin{figure}[htbp]
\centering
\colimg{M_chain.png}
\caption{决策链的六个环节。第 4 环把信息动作与品质动作分开（可分离），第 5 环只读硬测量（可拦截），第 1、2 环决定可重构，第 3 环决定可迁移}
\label{fig:chain}
\end{figure}

\begin{figure}[htbp]
\centering
\colimg{M_architecture.png}
\caption{五层闭环。左侧虚线为可安装观测进入估计层，右侧虚线为指令写回过程层；安全层位于链条末端，只读硬测量}
\label{fig:loop}
\end{figure}

实现中的时间步按这个顺序执行：读表面温度和平均含水，更新热阻代理和冷点估计，形成加热请求，用加热器温度削减请求，再推进十六条控制体。结构方程的路径积分和有限时域优化不在这个顺序里。热阻更新在是否更接近参照这一步使用了由真值温差构成的代理，因此方差下降必须连同该参照一起报告。顺序一旦固定，交叉表比较的对象也就固定：四条策略在同一积分器上的停机偏差，执行器接口始终只有四个量。

操作者需要的判断只有四个：能否停止，表面含水是否偏离目标，温度是否还在散开，继续施加边界是否安全。对应的隐藏量构成决策状态
\begin{equation}
\label{eq:z}
\begin{aligned}
\vect{z}(t)=\bigl(&
T_{\mathrm{cold}},\,
\bar T_s,\,
\sigma_T,\,
\bar\varepsilon_l,\,
\varepsilon_{l,s},\\
&\alpha_m,\,
\alpha_{b,s},\,
S_s,\,
m_{\mathrm{loss}}
\bigr).
\end{aligned}
\end{equation}
依次为冷点温度、表面均温、温度场标准差、平均液相分数、表面液相分数、致死进度、表面褐变、表面结壳与累计质量损失。停机看的是冷点和致死当量，不看全场的 \(L^2\) 范数。腔温只描述空气。新增一路观测必须降低上列至少一个判断的 \(\Delta\mathcal{I}\)，否则不进入观测向量。

主日志记录的只有 \(T_{\mathrm{cold}}\)、表面温度、平均含水与 \(\alpha_m\) 这几项，而九维决策状态里的其余分量并不落盘。九个分量中，停机真正读取的其实只有冷点与致死进度，表面液相和结壳则进入路径 B、C 的显示量。决策状态因此比日志宽、比全场窄：宽的那一侧用来说明还缺哪一路观测，窄的那一侧用来固定交叉表究竟比较哪几个量。食材类别只改先验与终端目标，分量定义本身保持不变——换一块土豆或一块面团，跟踪器仍读表面温度和平均含水，停机仍看 \(\alpha_m\)。

蒸、烤、组合不增加执行器，只限制四个量的允许集合：
\begin{equation}
\label{eq:mode}
(u,s,f,e)=
\begin{cases}
(\min(u,0.25),\,s,\,f,\,e) & \text{蒸}\\
(u,\,0,\,f,\,\max(e,0.45)) & \text{烤}\\
(u,\,s,\,f,\,e) & \text{组合}.
\end{cases}
\end{equation}
食品模块收到的是通量，不读取档位名称。面板上的三个字对应的是这三条不等式，不是三套独立的传热定律。蒸档把加热请求夹在 \(0.25\) 以下，蒸汽仍可离开零，冷凝潜热是这一档的主边界。烤档把蒸汽固定在零，并把排汽抬到至少 \(0.45\)，表面被推干，冷点只靠导热跟上。组合档放开四个量，蒸汽和干热可以在同一次过程里先后出现。

硬切断为
\begin{equation}
\label{eq:hard}
u_{\mathrm{cmd}}=
\begin{cases}
0, & \text{硬测量触发},\\
u_{\mathrm{req}}\,s(T_h), & T_h>230,\\
u_{\mathrm{req}}, & \text{其余},
\end{cases}
\end{equation}
其中 \(s(T_h)=1-0.75\,\mathrm{clip}\bigl((T_h-230)/40\bigr)\)。
加热器满功率平衡在 \(196.5^\circ\mathrm{C}\)，仿真中 \(230^\circ\mathrm{C}\) 削减不动作。规则仍与后验无关。水位干了仍加热，是干烧；门开着仍喷蒸汽，是烫伤；漏电仍通电，是电击。这三件事不随后验抖动，是因为硬切断的自变量里没有进度估计。

信息动作只占用 \([150,190)\) 这一段，其余时间的加热请求由估计温度、方差与厚度共同决定；安全层则插在请求与指令之间，只读硬测量。下文所说的“对偶”，指的正是这条固定窗加饱和线性请求的做法，“安全”指那道只读硬测量的投影。四条策略的差别集中在停止时刻怎么来：开环把它写成厚度的函数，腔温比例直接取时域终点，无探针与对偶则共用同一条饱和线性加热律。交叉表要比较的，就是这四条入口在同一积分器上给出的停机偏差。

整个过程里同时走着四个时间尺度，而它们按快慢排成一层一层的结构。毫秒级的剩余电流切断能够打断分钟级的食品估计，秒级的探针窗比内部热阻的时间常数还短，分钟级的品质加热又覆盖了剩下的大半过程，至于两次烹饪之间的先验写回，则发生在停机之后。这四层叠在一起，决定了安全层的时钟必须独立于估计器：快的那套保护不能被慢的那套计算拖住，漏电互感器也从不经过食品模型。家用电热器具上这件事早已成立，一体机要守住的，只是让进度估计同样走不进这条快回路。

\section{五条设计判据及其落地程度}

方案需要同时满足五条性质。只成立其中几条时，系统会退回到已经有的装置：只会定时，只会跟腔温，或只会在超温时断电。

\begin{enumerate}
\item \textbf{可重构。} 冷点温度、表面含水、反应进度和内部热阻能由有限观测与短时边界激励更新。稳态观测不含内部热阻，可重构依赖离开平衡的激励。
\item \textbf{可分离。} 同一次阀门动作要么为了降低后验熵，要么为了推动品质泛函。日志里要能查出这一时刻是哪一个函数发出的指令。
\item \textbf{可迁移。} 宏观传热、孔隙水分和反应进度耦合成低维决策状态。食材类别只改变先验和终端目标，不改变状态的分量定义。
\item \textbf{可更新。} 单次过程只更新事先声明的传递参数。结构方程的父节点不随这次数据改写。
\item \textbf{可拦截。} 品质模型的输出在进入执行器之前，被只依赖硬测量的不等式投影。硬保护不经过该模型。
\end{enumerate}

这五条都是设计判据，而不是事后总结。已经运行的一维核把“可分离”落实为一段 \(150\)--\(190\,\mathrm{s}\) 的固定窗：窗内蒸汽与风机按时间表离开平衡，窗外加热请求由冷点估计、方差和厚度决定。热阻更新用的参照由真值温差构成，估计器同时读取可安装通道与评分参照；要把可重构写严，第一步就是把参照换成可安装通道。安全投影则落在加热器温度这一路上：满功率平衡点 \(196.5^\circ\mathrm{C}\) 低于 \(230^\circ\mathrm{C}\) 阈值，削减在战役里保持待命，规则只读管温。

五条判据在后文的落地程度并不相同，表\ref{tab:properties}把它们逐条落到实现位置与验证证据上，表\ref{tab:properties}同时标出哪一条仍停在规格层：可更新是唯一没有日志支撑的一条，它需要把一次过程的后验写回下一次先验，而本文的战役结束于日志。

\begin{table*}[t]
\centering
\caption{五条设计判据的实现位置与验证证据。最后一行是唯一仍停在规格层的一条}
\label{tab:properties}
\begin{tabularx}{\textwidth}{l>{\raggedright\arraybackslash}p{3.6cm}>{\raggedright\arraybackslash}p{3.6cm}>{\raggedright\arraybackslash}p{3.2cm}}
\toprule
判据 & 实现位置 & 验证证据 & 状态 \\
\midrule
可重构 & 标量跟踪器与两状态 EKF；\(150\)--\(190\,\mathrm{s}\) 探针窗 & 表\ref{tab:est}：EKF 把偏差收窄、方差连续；S2 对偶 \(-0.112\) 对无探针 \(-0.579\) & 部分（参照仍用真值温差） \\
可分离 & 探针窗内蒸汽列跳到 \(1.0\)、窗外回基线；日志 \texttt{mode} 列 & 图\ref{fig:chain} 第 4 环；§5.4 轨迹与图\ref{fig:trpol} & 已实现 \\
可迁移 & 活度函数与食品类别只改先验与终端目标 & S1 三种食品的中位几乎重合（\(-0.338\) / \(-0.325\) / \(-0.313\)） & 已实现 \\
可更新 & 需把后验写回下一次先验 & 无日志支撑 & 仅规格 \\
可拦截 & 安全投影只读水位、门、漏电与管温 & 表\ref{tab:s2} 开门行；平衡点 \(196.5^\circ\mathrm{C}\) 低于阈值 & 已实现（削减待命） \\
\bottomrule
\end{tabularx}
\end{table*}

\section{本文的工作与结构}

\subsection{相对于已有工作，新增的是什么}

五条判据里，前四条在文献中都有各自的工具：逆热传导有反演理论，对偶控制有 Feldbaum 的经典结论，屏障函数有控制理论的形式保证，多孔介质的体积平均更是成熟框架。把这些工具找齐并不构成贡献。本文的出发点是：这些工具各自成立，却没有一处交代它们如何在同一台只有四个执行器的家用机器上接起来，以及接起来之后哪一条会先失效。

因此本文的贡献不在算法层，而在四件彼此可分开核对的事。

第一，把"内部热阻能否进入后验"写成一个可以在平衡点上判定真假的命题，并证明它在该点为假。两节点降阶方程中三个热阻与两个热容只以除数出现，对参数求导后恒为零；数值上三个不同的内部热阻在同一平衡点给出的冷点差小于 \(0.01\,\mathrm{K}\)。这解释了菜单时钟为什么必然失败——它闭合的是时间，而这一次过程里的热阻从未被激励——也把"加传感器有没有用"变成一个可以算的问题：稳态观测再密也不增加这一列的信息。

第二，给出一种只用四个执行器就能买到第二个灵敏度方向的激励方案，并说清它的边界。两段各 \(20\,\mathrm{s}\) 的探针先改相变符号、再改外部热阻，作用分别是把灵敏度矩阵从零秩恢复到满秩，以及把两个参数方向的张角拉大以压低条件数。这两件事常被混成一句"探针增加了信息"，但数学作用不同，失效条件也不同：前者要求离开平衡，后者要求窗口内的表面温差足够。

第三，把"信息不够"与"时间不够"分成两条可独立测量的轴。这是本文最实用的一条。厚件到不了设计线可以因为热阻没被看见，也可以因为看见了而扩散时间不足；在 \(55\,\mathrm{min}\) 上两者给出同一个观测，无法区分。把时域延到 \(70\,\mathrm{min}\) 之后方差不变而进度越过设计线，两者才分开。分开之后升级方向才唯一：加密探针没有用，延长时域或减薄才有用。

第四，用同一积分器上的 \(19440\) 与 \(11664\) 条样本给出四条策略的停机几何，并把"更好的估计器"与"更好的优化器"各自能买到什么分开报告。两状态滤波把热阻偏差从约 \(-3\) 收窄到 \(-1\) 左右、后验方差从两个取值变成连续分布，停机偏差却没有改善；滚动时域决策在三种厚度上都最省能量，却因共用有偏的冷点估计而更早停机。这条否定性结果之所以算贡献，是因为它可以被证伪：若瓶颈在优化，换 MPC 应同时改善两者，而实测只改善了能耗。

这四条分别落在第二到第五章，第六章把同一条状态流读回三维腔体，让链在时间上闭合。具体做过的工作如下。

\begin{enumerate}
\item 文献中的体积平均、品质动力学、逆热传导、对偶控制和屏障函数，分别标出继承层与规格层；商品路径与公开专利的信息止点整理成对照表，并画成一张四层矩阵。
\item 给出从多孔介质连续方程到十六条控制体、再到两节点工作方程的收缩，并定义九维决策状态；说明每一步收缩凭什么成立、代价是什么。
\item 证明平衡观测对内部热阻灵敏度为零；给出两段探针的设计，并更正它改善的是条件数而非秩；给出探针窗内冷点增量的上界。
\item 把同一次蒸汽分解为冷点潜热、表面蒸发和结壳热阻三条路径，安全投影的可行域与品质后验分离。
\item 在同一数值核上比较开环、腔温比例、无探针、对偶与滚动时域五种策略，并给出三档交叉、因子消融与网格收敛。
\item 实现读取同一状态流的三维腔体演示，边界、色带、路径与投影同屏可见，且演示读的是与主表同一次积分写出的数据。
\end{enumerate}

全书按问题推进，不按方法并列。第二章先划界：文献各自做实了哪一段，货架产品又闭合了哪一段，合起来标出这个缺口的位置。第三章回答“状态怎么来”，从三维连续方程一路收缩到十六条控制体，并说明每一步收缩凭什么成立。第四章回答“看得见吗”，先证明稳态观测对内部热阻的灵敏度为零，再给出两段探针这条出路，最后把安全投影钉在硬测量上。第五章回答“跑起来如何”，把同一积分器放到因子格子上，看厚度、初温和档位各自把结果推向哪里。第六章把同一条状态流读回三维腔体，让链在时间上闭合，再收束全文。

符号、离散格式、种子 7 的轨迹抽样和尺度估计放在附录 A 到 E；正文只保留论证需要的符号与表格，长表移到附录以便对照而不打断阅读。
